%% Idee cle : cas T grand (T=4, theta=0.5). Les impulsions temporelles
%% s'espacent, les coefficients Tc_n (plus nombreux, plus rapproches)
%% suivent de pres la meme enveloppe sinc. Repris du poly (section 1).
%% Consomme par slides.
\begin{tikzpicture}[x=0.7cm,y=0.7cm]
    \draw[-Latex] (-4.7,0) -- (4.7,0) node[right] {$t$};
    \draw[-Latex] (0,-0.2) -- (0,2.3) node[above] {};
    \foreach \c in {-4,0,4} {
        \draw[niceblue,thick] (\c-0.25,0) -- (\c-0.25,2) -- (\c+0.25,2) -- (\c+0.25,0);
    }
    \draw[latex-latex] (-0.25,2.5) -- (0.25,2.5) node[midway,above] {\footnotesize $\theta$};
    \node[left] at (0,2) {$\frac1\theta$};
    \node[below] at (4,0) {$T$};
\end{tikzpicture}
\hspace{0.5cm}
\begin{tikzpicture}
    \begin{axis}[
        width=6.5cm, height=4.3cm,
        xlabel={fréquence $f$}, ylabel={$T\,c_n$},
        xmin=-4, xmax=4, ymin=-0.3, ymax=1.1,
        axis lines=middle,
        xlabel style={at=(current axis.right of origin), anchor=west},
        ylabel style={at=(current axis.above origin), anchor=south},
    ]
    \addplot[domain=-4:4, samples=200, niceblue!40, thick] {sin(deg(1.570795*x))/(1.570795*x)};
    \addplot+[ycomb, black!70, thick, mark=*, mark size=1pt, mark options={fill=black!70}] coordinates {
        (-4.0000,-0.0000) (-3.7500,-0.0650) (-3.5000,-0.1286) (-3.2500,-0.1810) (-3.0000,-0.2122) (-2.7500,-0.2139) (-2.5000,-0.1801) (-2.2500,-0.1083) (-2.0000,0.0000) (-1.7500,0.1392) (-1.5000,0.3001) (-1.2500,0.4705) (-1.0000,0.6366) (-0.7500,0.7842) (-0.5000,0.9003) (-0.2500,0.9745) (0.0000,1.0000) (0.2500,0.9745) (0.5000,0.9003) (0.7500,0.7842) (1.0000,0.6366) (1.2500,0.4705) (1.5000,0.3001) (1.7500,0.1392) (2.0000,0.0000) (2.2500,-0.1083) (2.5000,-0.1801) (2.7500,-0.2139) (3.0000,-0.2122) (3.2500,-0.1810) (3.5000,-0.1286) (3.7500,-0.0650) (4.0000,-0.0000)
    };
    \end{axis}
\end{tikzpicture}
